\newpage
\subsection*{Abstract}

The interest in Bayesian designs that borrow historical information on the control arm is expected to increase following the
recent FDA guidance ``Use of Bayesian Methodology in Clinical Trials of Drug and Biological Products: Draft Guidance for Industry''.
Evaluating such designs is not straightforward, since borrowing breaks the frequentist guarantee of strict type I error (T1E) control. To address this,
an average T1E, which integrates the conditional T1E over a design prior for the control parameter, was recently proposed. This thesis
applies the same averaging to power, derives closed form expressions for both metrics under normal priors, obtains a closed form sample size formula
for a single normal analysis prior, and characterizes their sensitivity to misspecification of the design prior.
Prior-data conflict is handled by building the analysis prior as a robust meta-analytic-predictive (MAP) prior,
a mixture of a historical component and a weakly informative one.
The framework is then extended to dual criterion designs and to Bayesian group sequential designs. Designs are
compared through the Experimental Unit Information Index (EUII), a recently proposed per-patient metric
of evidentiary value which normalizes the diagnostic odds ratio
by the number of patients required. 
Because adding interim looks inflates the average T1E, an algorithm is proposed to calibrate designs to a common average T1E
and average power before the EUII is read. A Phase II Crohn's disease trial motivates and illustrates the comparisons throughout.

%%% Local Variables:
%%% mode: latex
%%% TeX-master: "MasterThesisSfS"
%%% End:
